\question{}

\begin{ابجد}
\item نامساوی \lr{Jensen} می‌گوید تابع $f: D \rightarrow \mathbb{R}$ روی دامنه‌ی $D$ محدب است اگر و تنها اگر:
\begin{equation*}
    \forall x, y \in D, \forall 0 \leq \theta \leq 1: f(\theta x + (1 - \theta)y) \leq \theta f(x) + (1 - \theta) f(y)
\end{equation*}
حال سعی می‌کنیم برای هر تابع داده شده نشان دهیم نامساوی جنسن برقرار است، یا مثال نقض بیاوریم:

\begin{itemize}
    \item $f(x) = \sqrt{x}$ بر روی $[0,\infty)$ \\ مقادیر $1, 9 \in [0,\infty)$ و $0 \leq t \leq 1$ را در نظر بگیرید:
          \begin{align*}
               & f((\frac{1}{2})1 + (1 - \frac{1}{2})9) = f(5) \simeq 2.24                            \\
               & \frac{1}{2}f(1) + (1 - \frac{1}{2})f(9) = 2                                          \\
               & f((\frac{1}{2})1 + (1 - \frac{1}{2})9) \nleq \frac{1}{2}f(1) + (1 - \frac{1}{2})f(9)
          \end{align*}
          مثال نقض آوردیم، پس نتیجه می‌گیریم که تابع $f(x)$ محدب نیست.

    \item $f(x) = x^2$ بر روی مجموعه‌ی کل اعداد حقیقی: \\ فرض می‌کنیم $0 \leq t \leq 1$ و $x$ و $y$ دو مقدار دلخواه باشند:
          \begin{align*}
              f(tx + (1 - t)y) \quad                & \leq \quad tf(x) + (1 - t)f(y)                               \\
              (tx + (1 - t)y)^2 \quad               & \leq \quad tx^2 + (1 - t)y^2                                 \\
              t^2x^2 + (1-t)^2y^2 + 2t(1-t)xy \quad & \leq \quad tx^2 + (1 - t)y^2                                 \\
              0 \quad                               & \leq \quad (t - t^2)x^2 + ((1-t) - (1-t)^2)y^2 - 2t(1 - t)xy \\
              0 \quad                               & \leq \quad t(1-t)(x^2 + y^2 - 2xy)                           \\
              0 \quad                               & \leq \quad t(1-t)(x - y)^2
          \end{align*}
          رابطه‌ی آخر همواره درست است. چون مقادیر $x$، $y$ و $t$ دلخواه بودند، به طور بازگشتی می‌توان نشان داد که نامساوی جنسن برای تمام مقادیر درست است. پس $f(x)$ محدب است.

\end{itemize}

\item می‌دانیم طبق آزمون مشتق دوم، اگر به ازای تمام نقاط دامنه $f'' (x) \geq 0$، تابع محدب است و اگر به ازای تمام نقاط دامنه $f''(x) \leq 0$، تابع مقعر است.

\begin{itemize}
    \item $f(x) = -\log(x)$ بر روی $(0, \infty)$ \\ فرض می‌کنیم لگاریتم همان $\ln$ باشد (در غیر این صورت همه چیز تنها در یک ضریب ثابت مثبت ضرب می‌شود):
          \begin{align*}
              f'(x) = -\frac{1}{x}, f''(x) = \frac{1}{x^2} \quad \Rightarrow \quad \forall x \in (0, \infty): f''(x) \geq 0
          \end{align*}
          بنابراین از آزمون مشتق دوم نتیجه می‌گیریم که تابع $f$ محدب است.

    \item $f(x) = x^4 - 2x^2$ بر روی مجموعه‌ی کل اعداد حقیقی \\
          \begin{align*}
              f'(x) = 4x^3 - 4x, f''(x) = 12x^2 - 4: \quad f''(0) = -4, f''(1) = 8
          \end{align*}
          دیدیم که $f''(x)$ می‌تواند هم مقادیر مثبت و هم مقادیر منفی بگیرد. به طور دقیق‌تر، تابع $f$ در بازه‌ی $(-\frac{1}{\sqrt{3}}, -\frac{1}{\sqrt{3}})$ مقعر و در خارج از این ناحیه محدب است.

\end{itemize}

\end{ابجد}